\documentclass{article}
\usepackage[T1]{fontenc}
\usepackage{lmodern}
\usepackage{graphicx}
\usepackage{amsmath,amssymb}
\usepackage[a4paper,margin=2cm]{geometry}
\usepackage[absolute,overlay]{textpos}
\setlength{\TPHorizModule}{1mm}
\setlength{\TPVertModule}{1mm}
\usepackage[indonesian]{babel}
\usepackage{hyperref}
\setlength{\emergencystretch}{2em}
% unit: Halaman 24

\begin{document}

% === Halaman 24 (python/native) ===

• Selanjutnya, kelompokkan menjadi kelompok-kelompok 7 bit:

0101010 1011100 1101100 1010010 0000011 0110101 1110010

0100000 0111001 1011001 0101100 0110111 0010011 0010101

1101000 0100000 0110101 1011001 0101111 001

• Konversi kelompok-kelompok 7-bit di atas menjadi nilai desimal

42 92 108 82 3 53 114 32 57 89 44 55 19 21 104 32 53 89 47 1.

24

\end{document}
